\documentclass[11pt]{article}
% =====================================================================
%  preamble.tex  —  shared preamble for Part I (physics.tex) and
%  Part II (math.tex) of "Bell Correlation as a Foundational Principle".
%
%  RECONSTRUCTED minimal preamble.  The original was referenced in
%  README_build_notes.md but never created as a file; physics/math were
%  only ever compiled against a throwaway stub.  This file collects the
%  macros and environments that physics.tex and math.tex are known to use
%  (from project records): a single shared theorem-numbering scheme across
%  both parts, the status labels, the NSA symbols (st, ord, *R), and the
%  categorical chain C_corr -> C_dist -> C_act -> C_eq.
%
%  BUILD (per README_build_notes.md), bidirectional xr cross-references:
%    each of physics.tex / math.tex does % =====================================================================
%  preamble.tex  —  shared preamble for Part I (physics.tex) and
%  Part II (math.tex) of "Bell Correlation as a Foundational Principle".
%
%  RECONSTRUCTED minimal preamble.  The original was referenced in
%  README_build_notes.md but never created as a file; physics/math were
%  only ever compiled against a throwaway stub.  This file collects the
%  macros and environments that physics.tex and math.tex are known to use
%  (from project records): a single shared theorem-numbering scheme across
%  both parts, the status labels, the NSA symbols (st, ord, *R), and the
%  categorical chain C_corr -> C_dist -> C_act -> C_eq.
%
%  BUILD (per README_build_notes.md), bidirectional xr cross-references:
%    each of physics.tex / math.tex does % =====================================================================
%  preamble.tex  —  shared preamble for Part I (physics.tex) and
%  Part II (math.tex) of "Bell Correlation as a Foundational Principle".
%
%  RECONSTRUCTED minimal preamble.  The original was referenced in
%  README_build_notes.md but never created as a file; physics/math were
%  only ever compiled against a throwaway stub.  This file collects the
%  macros and environments that physics.tex and math.tex are known to use
%  (from project records): a single shared theorem-numbering scheme across
%  both parts, the status labels, the NSA symbols (st, ord, *R), and the
%  categorical chain C_corr -> C_dist -> C_act -> C_eq.
%
%  BUILD (per README_build_notes.md), bidirectional xr cross-references:
%    each of physics.tex / math.tex does \input{preamble} then
%    \usepackage{xr}\externaldocument[I-]{physics}  (in math.tex)
%    \usepackage{xr}\externaldocument[II-]{math}     (in physics.tex)
%    Compile order: bootstrap each once (pdflatex) to make .aux, then
%    latexmk both twice to converge cross-document refs.
%
%  This preamble is deliberately conservative: it defines every macro the
%  two documents are known to reference, plus the common analysis symbols,
%  so a missing \input does not break the build.  If physics/math use a
%  macro not defined here, add it here (this is the single shared source).
%  files own \documentclass (verified: both real files issue
%  \documentclass[11pt]{article} on line 1), so this preamble must NOT.
% =====================================================================

% ---- encoding / fonts / layout ----
\usepackage[T1]{fontenc}
\usepackage[utf8]{inputenc}
\usepackage[margin=1.1in]{geometry}

% ---- math ----
\usepackage{amsmath,amsthm,amssymb,mathtools}
\usepackage{bm}

% ---- graphics / tables / misc ----
\usepackage{graphicx}
\usepackage{booktabs}
\usepackage{enumitem}
\usepackage{xcolor}

% ---- hyperref last (before xr in the document) ----
\usepackage[colorlinks=true,linkcolor=blue,citecolor=blue,urlcolor=blue]{hyperref}

% =====================================================================
%  Theorem environments — ONE shared numbering scheme across both parts.
%  All share the counter "thmcounter" numbered within section, so Part I
%  and Part II do not collide when cross-referenced.
% =====================================================================
\newtheorem{theorem}{Theorem}[section]
\newtheorem{proposition}[theorem]{Proposition}
\newtheorem{lemma}[theorem]{Lemma}
\newtheorem{corollary}[theorem]{Corollary}
\theoremstyle{definition}
\newtheorem{definition}[theorem]{Definition}
\newtheorem{example}[theorem]{Example}
\theoremstyle{remark}
\newtheorem{remark}[theorem]{Remark}
\newtheorem{note}[theorem]{Note}
\newtheorem{principle}[theorem]{Principle}

% =====================================================================
%  Status labels (README §7: TeX holds only established statements;
%  these tag the epistemic status inline).
% =====================================================================
\newcommand{\established}{\textcolor{green!55!black}{\textsf{[established]}}}
\newcommand{\conditional}{\textcolor{orange!85!black}{\textsf{[conditional]}}}
\newcommand{\deferred}{\textcolor{blue!70!black}{\textsf{[deferred]}}}
\newcommand{\TODO}{\textcolor{red!80!black}{\textsf{[TODO]}}}

% =====================================================================
%  Nonstandard-analysis symbols.
%  *R = hyperreals; st = standard-part map; ord = infinitesimal order;
%  fin = finite hyperreals. eps is the fixed positive infinitesimal
%  lattice spacing (with eps^2 > 0 still infinitesimal).
% =====================================================================
\newcommand{\stR}{{}^{*}\mathbb{R}}        % hyperreals
\newcommand{\stN}{{}^{*}\mathbb{N}}        % hypernaturals
\DeclareMathOperator{\st}{st}              % standard-part map
\DeclareMathOperator{\ord}{ord}            % infinitesimal order
\DeclareMathOperator{\fin}{fin}            % finite hyperreals
\newcommand{\eps}{\varepsilon}

% measurement map  M(f) = st( f / eps^{ord f} )
\newcommand{\Meas}{\mathcal{M}}

% =====================================================================
%  Categorical chain  C_corr -> C_dist -> C_act -> C_eq.
% =====================================================================
\newcommand{\Ccorr}{\mathcal{C}_{\mathrm{corr}}}
\newcommand{\Cdist}{\mathcal{C}_{\mathrm{dist}}}
\newcommand{\Cact}{\mathcal{C}_{\mathrm{act}}}
\newcommand{\Ceq}{\mathcal{C}_{\mathrm{eq}}}

% =====================================================================
%  Common shorthands used across the two parts.
% =====================================================================
\newcommand{\R}{\mathbb{R}}
\newcommand{\N}{\mathbb{N}}
\newcommand{\Z}{\mathbb{Z}}
\newcommand{\C}{\mathbb{C}}
\newcommand{\Q}{\mathbb{Q}}
\newcommand{\corr}{C}                      % correlation C(theta) = -cos theta
\newcommand{\dist}{D}                      % relational distance D = -log|C|
\newcommand{\elliptope}{\mathcal{E}}
\newcommand{\Gram}{G}
\newcommand{\inner}[2]{\langle #1, #2 \rangle}
\newcommand{\abs}[1]{\left| #1 \right|}
\newcommand{\norm}[1]{\left\lVert #1 \right\rVert}
\newcommand{\set}[1]{\left\{ #1 \right\}}
\DeclareMathOperator{\rank}{rank}
\DeclareMathOperator{\tr}{tr}
\DeclareMathOperator{\diag}{diag}
\DeclareMathOperator*{\argmin}{arg\,min}
\DeclareMathOperator*{\argmax}{arg\,max}

% Bell-framework specifics: half-angle delta, Boole quantities F_eps, the
% window depth m = min_eps F_eps.
\newcommand{\Fbool}{F}
\newcommand{\windowdepth}{m}

% =====================================================================
%  Macros required by the real physics.tex / math.tex (verified by macro
%  extraction). Meanings inferred from usage context.
% =====================================================================
\newcommand{\RR}{\mathbb{R}}               % reals (used as \RR)
\newcommand{\ZZ}{\mathbb{Z}}               % integers (Z = pi_1(S^1))
\newcommand{\Cmat}{C}                       % correlation matrix / value, used as |\Cmat|
\newcommand{\dent}{d}                       % relational distance d=-log|C|
\newcommand{\Tcirc}{\mathbb{T}}            % torus; used as \Tcirc^{3}, \Tcirc^{p}
\newcommand{\Scirc}{\mathbb{S}}            % circle S^1; Z \cong pi_1(\Scirc)
\newcommand{\LamQCD}{\Lambda_{\mathrm{QCD}}}% QCD scale
\newcommand{\Eang}{E}                       % E_ang: correlation E(theta) as a function of angle; \Eang(theta)=cos 2theta/(2N)

% end of preamble.tex
 then
%    \usepackage{xr}\externaldocument[I-]{physics}  (in math.tex)
%    \usepackage{xr}\externaldocument[II-]{math}     (in physics.tex)
%    Compile order: bootstrap each once (pdflatex) to make .aux, then
%    latexmk both twice to converge cross-document refs.
%
%  This preamble is deliberately conservative: it defines every macro the
%  two documents are known to reference, plus the common analysis symbols,
%  so a missing \input does not break the build.  If physics/math use a
%  macro not defined here, add it here (this is the single shared source).
%  files own \documentclass (verified: both real files issue
%  \documentclass[11pt]{article} on line 1), so this preamble must NOT.
% =====================================================================

% ---- encoding / fonts / layout ----
\usepackage[T1]{fontenc}
\usepackage[utf8]{inputenc}
\usepackage[margin=1.1in]{geometry}

% ---- math ----
\usepackage{amsmath,amsthm,amssymb,mathtools}
\usepackage{bm}

% ---- graphics / tables / misc ----
\usepackage{graphicx}
\usepackage{booktabs}
\usepackage{enumitem}
\usepackage{xcolor}

% ---- hyperref last (before xr in the document) ----
\usepackage[colorlinks=true,linkcolor=blue,citecolor=blue,urlcolor=blue]{hyperref}

% =====================================================================
%  Theorem environments — ONE shared numbering scheme across both parts.
%  All share the counter "thmcounter" numbered within section, so Part I
%  and Part II do not collide when cross-referenced.
% =====================================================================
\newtheorem{theorem}{Theorem}[section]
\newtheorem{proposition}[theorem]{Proposition}
\newtheorem{lemma}[theorem]{Lemma}
\newtheorem{corollary}[theorem]{Corollary}
\theoremstyle{definition}
\newtheorem{definition}[theorem]{Definition}
\newtheorem{example}[theorem]{Example}
\theoremstyle{remark}
\newtheorem{remark}[theorem]{Remark}
\newtheorem{note}[theorem]{Note}
\newtheorem{principle}[theorem]{Principle}

% =====================================================================
%  Status labels (README §7: TeX holds only established statements;
%  these tag the epistemic status inline).
% =====================================================================
\newcommand{\established}{\textcolor{green!55!black}{\textsf{[established]}}}
\newcommand{\conditional}{\textcolor{orange!85!black}{\textsf{[conditional]}}}
\newcommand{\deferred}{\textcolor{blue!70!black}{\textsf{[deferred]}}}
\newcommand{\TODO}{\textcolor{red!80!black}{\textsf{[TODO]}}}

% =====================================================================
%  Nonstandard-analysis symbols.
%  *R = hyperreals; st = standard-part map; ord = infinitesimal order;
%  fin = finite hyperreals. eps is the fixed positive infinitesimal
%  lattice spacing (with eps^2 > 0 still infinitesimal).
% =====================================================================
\newcommand{\stR}{{}^{*}\mathbb{R}}        % hyperreals
\newcommand{\stN}{{}^{*}\mathbb{N}}        % hypernaturals
\DeclareMathOperator{\st}{st}              % standard-part map
\DeclareMathOperator{\ord}{ord}            % infinitesimal order
\DeclareMathOperator{\fin}{fin}            % finite hyperreals
\newcommand{\eps}{\varepsilon}

% measurement map  M(f) = st( f / eps^{ord f} )
\newcommand{\Meas}{\mathcal{M}}

% =====================================================================
%  Categorical chain  C_corr -> C_dist -> C_act -> C_eq.
% =====================================================================
\newcommand{\Ccorr}{\mathcal{C}_{\mathrm{corr}}}
\newcommand{\Cdist}{\mathcal{C}_{\mathrm{dist}}}
\newcommand{\Cact}{\mathcal{C}_{\mathrm{act}}}
\newcommand{\Ceq}{\mathcal{C}_{\mathrm{eq}}}

% =====================================================================
%  Common shorthands used across the two parts.
% =====================================================================
\newcommand{\R}{\mathbb{R}}
\newcommand{\N}{\mathbb{N}}
\newcommand{\Z}{\mathbb{Z}}
\newcommand{\C}{\mathbb{C}}
\newcommand{\Q}{\mathbb{Q}}
\newcommand{\corr}{C}                      % correlation C(theta) = -cos theta
\newcommand{\dist}{D}                      % relational distance D = -log|C|
\newcommand{\elliptope}{\mathcal{E}}
\newcommand{\Gram}{G}
\newcommand{\inner}[2]{\langle #1, #2 \rangle}
\newcommand{\abs}[1]{\left| #1 \right|}
\newcommand{\norm}[1]{\left\lVert #1 \right\rVert}
\newcommand{\set}[1]{\left\{ #1 \right\}}
\DeclareMathOperator{\rank}{rank}
\DeclareMathOperator{\tr}{tr}
\DeclareMathOperator{\diag}{diag}
\DeclareMathOperator*{\argmin}{arg\,min}
\DeclareMathOperator*{\argmax}{arg\,max}

% Bell-framework specifics: half-angle delta, Boole quantities F_eps, the
% window depth m = min_eps F_eps.
\newcommand{\Fbool}{F}
\newcommand{\windowdepth}{m}

% =====================================================================
%  Macros required by the real physics.tex / math.tex (verified by macro
%  extraction). Meanings inferred from usage context.
% =====================================================================
\newcommand{\RR}{\mathbb{R}}               % reals (used as \RR)
\newcommand{\ZZ}{\mathbb{Z}}               % integers (Z = pi_1(S^1))
\newcommand{\Cmat}{C}                       % correlation matrix / value, used as |\Cmat|
\newcommand{\dent}{d}                       % relational distance d=-log|C|
\newcommand{\Tcirc}{\mathbb{T}}            % torus; used as \Tcirc^{3}, \Tcirc^{p}
\newcommand{\Scirc}{\mathbb{S}}            % circle S^1; Z \cong pi_1(\Scirc)
\newcommand{\LamQCD}{\Lambda_{\mathrm{QCD}}}% QCD scale
\newcommand{\Eang}{E}                       % E_ang: correlation E(theta) as a function of angle; \Eang(theta)=cos 2theta/(2N)

% end of preamble.tex
 then
%    \usepackage{xr}\externaldocument[I-]{physics}  (in math.tex)
%    \usepackage{xr}\externaldocument[II-]{math}     (in physics.tex)
%    Compile order: bootstrap each once (pdflatex) to make .aux, then
%    latexmk both twice to converge cross-document refs.
%
%  This preamble is deliberately conservative: it defines every macro the
%  two documents are known to reference, plus the common analysis symbols,
%  so a missing \input does not break the build.  If physics/math use a
%  macro not defined here, add it here (this is the single shared source).
%  files own \documentclass (verified: both real files issue
%  \documentclass[11pt]{article} on line 1), so this preamble must NOT.
% =====================================================================

% ---- encoding / fonts / layout ----
\usepackage[T1]{fontenc}
\usepackage[utf8]{inputenc}
\usepackage[margin=1.1in]{geometry}

% ---- math ----
\usepackage{amsmath,amsthm,amssymb,mathtools}
\usepackage{bm}

% ---- graphics / tables / misc ----
\usepackage{graphicx}
\usepackage{booktabs}
\usepackage{enumitem}
\usepackage{xcolor}

% ---- hyperref last (before xr in the document) ----
\usepackage[colorlinks=true,linkcolor=blue,citecolor=blue,urlcolor=blue]{hyperref}

% =====================================================================
%  Theorem environments — ONE shared numbering scheme across both parts.
%  All share the counter "thmcounter" numbered within section, so Part I
%  and Part II do not collide when cross-referenced.
% =====================================================================
\newtheorem{theorem}{Theorem}[section]
\newtheorem{proposition}[theorem]{Proposition}
\newtheorem{lemma}[theorem]{Lemma}
\newtheorem{corollary}[theorem]{Corollary}
\theoremstyle{definition}
\newtheorem{definition}[theorem]{Definition}
\newtheorem{example}[theorem]{Example}
\theoremstyle{remark}
\newtheorem{remark}[theorem]{Remark}
\newtheorem{note}[theorem]{Note}
\newtheorem{principle}[theorem]{Principle}

% =====================================================================
%  Status labels (README §7: TeX holds only established statements;
%  these tag the epistemic status inline).
% =====================================================================
\newcommand{\established}{\textcolor{green!55!black}{\textsf{[established]}}}
\newcommand{\conditional}{\textcolor{orange!85!black}{\textsf{[conditional]}}}
\newcommand{\deferred}{\textcolor{blue!70!black}{\textsf{[deferred]}}}
\newcommand{\TODO}{\textcolor{red!80!black}{\textsf{[TODO]}}}

% =====================================================================
%  Nonstandard-analysis symbols.
%  *R = hyperreals; st = standard-part map; ord = infinitesimal order;
%  fin = finite hyperreals. eps is the fixed positive infinitesimal
%  lattice spacing (with eps^2 > 0 still infinitesimal).
% =====================================================================
\newcommand{\stR}{{}^{*}\mathbb{R}}        % hyperreals
\newcommand{\stN}{{}^{*}\mathbb{N}}        % hypernaturals
\DeclareMathOperator{\st}{st}              % standard-part map
\DeclareMathOperator{\ord}{ord}            % infinitesimal order
\DeclareMathOperator{\fin}{fin}            % finite hyperreals
\newcommand{\eps}{\varepsilon}

% measurement map  M(f) = st( f / eps^{ord f} )
\newcommand{\Meas}{\mathcal{M}}

% =====================================================================
%  Categorical chain  C_corr -> C_dist -> C_act -> C_eq.
% =====================================================================
\newcommand{\Ccorr}{\mathcal{C}_{\mathrm{corr}}}
\newcommand{\Cdist}{\mathcal{C}_{\mathrm{dist}}}
\newcommand{\Cact}{\mathcal{C}_{\mathrm{act}}}
\newcommand{\Ceq}{\mathcal{C}_{\mathrm{eq}}}

% =====================================================================
%  Common shorthands used across the two parts.
% =====================================================================
\newcommand{\R}{\mathbb{R}}
\newcommand{\N}{\mathbb{N}}
\newcommand{\Z}{\mathbb{Z}}
\newcommand{\C}{\mathbb{C}}
\newcommand{\Q}{\mathbb{Q}}
\newcommand{\corr}{C}                      % correlation C(theta) = -cos theta
\newcommand{\dist}{D}                      % relational distance D = -log|C|
\newcommand{\elliptope}{\mathcal{E}}
\newcommand{\Gram}{G}
\newcommand{\inner}[2]{\langle #1, #2 \rangle}
\newcommand{\abs}[1]{\left| #1 \right|}
\newcommand{\norm}[1]{\left\lVert #1 \right\rVert}
\newcommand{\set}[1]{\left\{ #1 \right\}}
\DeclareMathOperator{\rank}{rank}
\DeclareMathOperator{\tr}{tr}
\DeclareMathOperator{\diag}{diag}
\DeclareMathOperator*{\argmin}{arg\,min}
\DeclareMathOperator*{\argmax}{arg\,max}

% Bell-framework specifics: half-angle delta, Boole quantities F_eps, the
% window depth m = min_eps F_eps.
\newcommand{\Fbool}{F}
\newcommand{\windowdepth}{m}

% =====================================================================
%  Macros required by the real physics.tex / math.tex (verified by macro
%  extraction). Meanings inferred from usage context.
% =====================================================================
\newcommand{\RR}{\mathbb{R}}               % reals (used as \RR)
\newcommand{\ZZ}{\mathbb{Z}}               % integers (Z = pi_1(S^1))
\newcommand{\Cmat}{C}                       % correlation matrix / value, used as |\Cmat|
\newcommand{\dent}{d}                       % relational distance d=-log|C|
\newcommand{\Tcirc}{\mathbb{T}}            % torus; used as \Tcirc^{3}, \Tcirc^{p}
\newcommand{\Scirc}{\mathbb{S}}            % circle S^1; Z \cong pi_1(\Scirc)
\newcommand{\LamQCD}{\Lambda_{\mathrm{QCD}}}% QCD scale
\newcommand{\Eang}{E}                       % E_ang: correlation E(theta) as a function of angle; \Eang(theta)=cos 2theta/(2N)

% end of preamble.tex

\usepackage{xr}
\externaldocument[I-]{physics}

\title{Bell Correlation as a Foundational Principle\\
\large Part II --- Mathematics: Standard-Part Measurement, Spectral Geometry, and the Mass Gap}
\author{T. Yanagisawa}
\date{\today}

\begin{document}
\maketitle

\begin{abstract}
This part develops the mathematical backbone behind the physics reconstruction
(Part~I): the standard-part map $\st$ as the universal model of measurement via
non-standard analysis, the spectral/MDS recovery of geometry from the correlation
matrix, and the mass-gap argument from the Bell-derived transfer operator. The
positivity of the gap is the mathematical result; the value $\LamQCD$ enters as an
empirical input. The transfer-operator argument is presented under the explicit
hypothesis~(H1) identifying the Bell two-point function with the Euclidean
Yang--Mills two-point function.
\end{abstract}

\tableofcontents

% ======================================================================
\section{Related work and positioning}\label{sec:math-related}
The use of nonstandard analysis and the standard-part map in mathematical physics is
established: hyperfinite and hypercontinuous methods~\cite{albeverio}, Yamashita's
hyperfinite formalism of quantum fields~\cite{yamashita}, and models employing a
hyperfinite real axis with infinitesimal lattice spacing as the support of physical
phenomena~\cite{nsa-qm}. Our graded measurement $M(f)=\st(f/\varepsilon^{\operatorname{ord}(f)})$,
with distinct infinitesimal orders for space and time ($\varepsilon^2>0$), follows this
lineage; the novelty is its use as the \emph{universal model of Bell measurement} feeding
the correlation-to-geometry construction, and the observation (\S\ref{sec:st}) that $\st$
constrains neither the spatial dimension nor the topology.

On the relativistic-causality side, Finberg~\cite{finberg} argues, within locally
covariant AQFT, that Bell local causality is incompatible with diffeomorphism
covariance---any screening-off completion requires preferred foliations or non-covariant
beables. That result \emph{assumes} spacetime and shows Bell correlations resist a
covariant classical-causal embedding; our work runs in the opposite direction, taking
Bell correlations as primary and \emph{outputting} a spacetime whose preferred time
arises from covering~II. The two are complementary: Finberg's preferred-foliation
requirement is the time-asymmetry that covering~II supplies constructively.

% ======================================================================
\section{Measurement as the standard-part map}\label{sec:st}
\established\ (foundational, README \S2c)

\begin{definition}[Graded measurement]\label{def:st}
Let $\varepsilon>0$ be a positive infinitesimal lattice spacing (in the nonstandard sense,
with $\varepsilon^{2}>0$ still infinitesimal). For a nonstandard quantity $f$ of
infinitesimal order $\operatorname{ord}(f)$, the \emph{graded measurement} is
\begin{equation}
M(f)=\st\!\left(\frac{f}{\varepsilon^{\operatorname{ord}(f)}}\right),
\label{eq:graded}
\end{equation}
where $\st$ is the standard-part map sending a limited hyperreal to the unique real
infinitely close to it. Bell measurement is one instance of this single operation.
\end{definition}

\begin{proposition}[Automatic exclusion of non-near-standard quantities]\label{prop:exclude}
At settings where the correlation vanishes ($\cos2\theta=0$), the entanglement distance
$\dent=-\log|\Cmat|$ is unlimited and therefore has no standard part; such points are
automatically excluded from the measured geometry. The excluded set has measure zero.
\end{proposition}

\begin{proof}
If $|\Cmat|\to0$ then $-\log|\Cmat|\to+\infty$, an unlimited hyperreal, which lies outside
the domain of $\st$ and so is not assigned any real value by~\eqref{eq:graded}. The
vanishing of $\cos2\theta$ occurs on a discrete set of angles, of Lebesgue measure zero in
the setting space.
\end{proof}

\begin{remark}[$\st$ fixes neither dimension nor topology]\label{rem:st-limits}
The standard-part map constrains neither the number of independent settings $p$ nor the
choice of covering. It is definable identically whether the support is $\Scirc$ or its
cover $\RR$, and places no bound on $p$ (notebook~17, \S\S17.9--17.11). Hence the spatial
dimension and the time topology lie \emph{upstream} of $\st$; this is the measurement-side
counterpart of the external-input status of dimension three established in Part~I,
\S\ref{I-sec:dim}.
\end{remark}

\begin{remark}[$\st$ is algebraically neutral]\label{rem:st-algebra}
The standard-part map likewise constrains no algebraic structure of the observables. For
limited-hyperreal-valued operators on a hyperfinite-dimensional ${}^{*}$Hilbert space (in the
sense of~\cite{albeverio,yamashita}), $\st$ acts entrywise and commutes with products and
commutators, $\st([A,B])=[\st A,\st B]$; a noncommuting pair remains noncommuting after $\st$
and a commuting pair remains commuting (notebook~33, \S33.1, verified symbolically with an
infinitesimal perturbation). Thus $\st$ neither creates nor destroys noncommutativity: just as
it fixes neither dimension nor topology (Rmk.~\ref{rem:st-limits}), it fixes no commutation
relations. The projective/noncommutative structure of measurement therefore lies neither
downstream nor upstream of $\st$, but is an independent datum (\S\ref{sec:qcorr}).
\end{remark}

\begin{remark}[$\st$ grades modes but does not order them]\label{rem:st-ordering}
A third structure lies upstream of the graded measurement, surfacing when the correlation
is resolved into spectral modes. Write a stationary correlation in spectral form with each
mode carried at a distinct infinitesimal order,
$\Cmat(t)=c_0+\sum_{k\ge1}\varepsilon^{k}c_k\cos(\omega_k t)$; then the graded
measurement~\eqref{eq:graded} at order $k$ returns exactly the $k$-th mode, so $M$ reads out
one mode per order and $\operatorname{ord}$ supplies an \emph{integer} ordering of the modes.
The rank of the correlation, however, only \emph{counts} modes: it is invariant under
relabelling which frequency is called mode~$1,2,\dots$ (notebooks~35--36, verified
numerically). Hence the assignment of a given physical quantity to a given order---the datum
that turns the mode count into an ordered tower---is not fixed by $\st$ or by the rank; it is
a modelling input. Just as $\st$ fixes neither dimension nor topology
(Rmk.~\ref{rem:st-limits}) nor the commutation relations (Rmk.~\ref{rem:st-algebra}), it does
not fix the mode ordering: the order assignment is an external input. This is the
measurement-side locus of the conditional bridge between the rank-two rigidity of
\S\ref{sec:nophase} and the minimal-zero construction of the companion
paper~\cite{rank2paper} (the identification ``minimal zero $=$ graded standard part'' holds
\emph{given} the order assignment, not unconditionally), and it assigns $\varepsilon$ no
numeric scale, consistent with the ordinal role of $\operatorname{ord}$ in~\eqref{eq:graded}.
\end{remark}

% ======================================================================
\section{Spectral recovery of geometry (MDS)}\label{sec:mds}
\established\ (notebooks 01, 04, 16, 18)

\begin{definition}[Correlation matrix and centred Gram]\label{def:gram}
Given pairwise distances $D_{ij}=\dent(\theta_i,\theta_j)$, form the squared-distance
matrix $D^{(2)}$ with entries $D_{ij}^2$ and the centred Gram matrix
\begin{equation}
B=-\tfrac12\,J\,D^{(2)}\,J,\qquad J=I-\tfrac1m\mathbf{1}\mathbf{1}^{\top},
\label{eq:gram}
\end{equation}
where $m$ is the number of settings. The positive eigenvalues of $B$ and their
eigenvectors give the MDS coordinates of the emergent geometry.
\end{definition}

\begin{proposition}[Additive decomposition of distance]\label{prop:additive}
For $p$ independent angular settings, the setting space is the torus $\Tcirc^{p}$ and the
squared distance decomposes additively,
$\dent^2=\sum_{c=1}^{p}\dent_c^{2}$, over the $p$ axes. The intrinsic dimension of the MDS
image is $p$ and its embedding dimension is $2p$.
\end{proposition}

\begin{proof}
Independent settings contribute independent angular factors, each an $\Scirc$ with its own
$\cos2$ correlation, so distances add in quadrature (notebook~4). Each $\Scirc$ requires
two MDS coordinates (Part~I, Lemma~\ref{I-lem:rank2}) but is intrinsically one-dimensional;
over $p$ axes this gives embedding dimension $2p$ and intrinsic dimension $p$.
\end{proof}

\begin{proposition}[Recovery of $\Tcirc^{3}$]\label{prop:t3}
For $p=3$, MDS of $|\Cmat|$ recovers an image of embedding dimension six and intrinsic
dimension three, with the three axis subspaces mutually orthogonal.
\end{proposition}

\begin{proof}
By Proposition~\ref{prop:additive} with $p=3$, the leading spectrum of $B$ consists of six
equal dominant eigenvalues (three axes $\times$ two coordinates) with a sharp drop to the
seventh, identifying intrinsic dimension three. Computing the principal angles between the
two-dimensional subspaces of distinct axes yields zero overlap, confirming orthogonality
(notebook~18, part~a).
\end{proof}

% ======================================================================
\section{No phase transition: rank-two rigidity}\label{sec:nophase}
\established\ (notebook 20b; supersedes the inconclusive notebook 20)

A natural question is whether the emergent geometry undergoes a phase transition---a
geometrogenesis---as some control parameter varies, in the spirit of emergent-geometry
studies across quantum phases~\cite{emergentdist2025,persistenthom}. We answer in the
negative for the given correlation, and locate the reason in the rigidity of its rank.

\begin{proposition}[Rank-two rigidity]\label{prop:nophase}
Built coordinate-free from $\Cmat_{ij}=\cos2(\theta_i-\theta_j)/(2N)$ via
$\dent=-\log|\Cmat|$ and MDS, the correlation matrix $\Cmat$ has rank two (the circle of
Part~I, Lemma~\ref{I-lem:bounded}) independently of $N$, of the number of settings $L$,
and in the continuum limit $\varepsilon\to0$. Consequently no order parameter exhibits
non-analytic behaviour: there is no phase transition.
\end{proposition}

\begin{proof}
On a near-standard window $|\Delta\theta|<\pi/4$ the distances are finite and the MDS is
well defined; the eigenvalues of $\Cmat$ show a constant rank two with $\lambda_3/\lambda_1
\to0$ as $L\to\infty$ (notebook~20b, verified to $\sim10^{-16}$). The order-parameter
candidates (largest MDS eigenvalue, spectral gap, effective dimension) vary smoothly as the
window approaches the zero $\Delta\theta=\pi/4$; no jump occurs. The earlier notebook~20
reported ``no transition'' for a different and flawed reason---it assumed lattice
coordinates and an exponential kernel $e^{-d/\xi}$ rather than the given $\cos2\theta$, so
its scaling law was a tautology in which $\xi$ was a mere overall factor. The present
result is coordinate-free and rests on the structural rigidity of the rank, not on a
tautology. The zeros of $\cos2\theta$ ($\Delta\theta=\pi/4+n\pi/2$, where $\Cmat=0$ and
$\dent=\infty$) produce excluded pairs---a non-standard, measure-zero set excised by the
standard part (Part~II analogue of the exclusion in \S\ref{sec:st})---which mark the edge
of the near-standard window, not a phase boundary. NSA: rank two is an internal property
preserved under $\st$ across the hyperfinite refinement.
\end{proof}

\begin{remark}[Relation to mutual-information emergent distance]\label{rem:nophase}
The mutual-information construction $d_E\propto1/\sqrt I$
of~\cite{emergentdist2025,persistenthom} ties the \emph{metricity} of the emergent distance
to the decay law of the correlation: a power law (critical regime, no characteristic
length) gives a valid metric, exponential decay (gapped regime) violates the triangle
inequality. The given $\cos2\theta$ neither decays nor has a characteristic length; in that
classification it sits on the metric (critical-like) side, and indeed yields the stable
rank-two geometry above. There is no control parameter---no analogue of the XXZ
anisotropy---to separate phases, so frameworks that read spacetime emergence as a phase
transition (e.g.\ group-field-theory condensation) have no point of application here. This
is a clean difference of setting, not a disagreement of result.
\end{remark}

\begin{remark}[The rigidity as level one of a graded tower]\label{rem:tower}
\conditional\ (companion rank-two paper, \S5.5--5.6; notebooks 34--36)
The rank-two rigidity of Prop.~\ref{prop:nophase} is the first level of a
graded structure analysed in the companion paper on the rank-two edge of
the correlation elliptope~\cite{rank2paper}. There the inverse question is
posed in the manner of the amplitudes result of Cheung et
al.~\cite{cheung2025}: a stationary correlation kernel whose every
$3\times3$ Gram matrix is rank-deficient (the mandated zero) and is of rank
\emph{exactly} two (minimality), together with continuity, forces
$\Cmat(t)=\cos(\omega t)$; raising the number of settings $L$ and demanding
rank exactly $L-1$ grades this into a tower, with
$\operatorname{rank}=2K+[\text{constant mode}]$ for $K$ spectral modes, so
that rank $(L-1)$ leaves no room for an independent order-$L$ correlation.
The rank-two case here is level one: a single harmonic, no room for an
independent three-point term. This is a structural---not dynamical---reason
for the absence of a phase transition: a rank-two stationary kernel cannot
cross to a qualitatively different regime without leaving the rank-two
locus.
\end{remark}

\begin{remark}[Minimal zero as a graded standard part]\label{rem:tower-nsa}
\conditional\ (notebooks 35--36)
The tower meshes with the graded measurement of \S\ref{sec:st}. Writing a
correlation in spectral form with each mode carried at a distinct
infinitesimal order, $\Cmat(t)=c_0+\sum_k\varepsilon^{k}c_k\cos(\omega_k t)$,
the graded measurement~\eqref{eq:graded} at order $k$ returns exactly the
$k$-th mode, $M$ reading out one mode per order. The minimal-zero cutoff
that fixes the rank then coincides with the standard-part truncation that
discards orders above the measured one: \emph{minimal zero $=$ graded
standard part}, with $\operatorname{ord}$ supplying the integer ordering of
modes that the rank alone does not (rank counts modes but is blind to their
labelling). The identification is conditional on the modelling choice of
which quantity sits at which order---it is not forced by rank---and assigns
$\varepsilon$ no numeric scale, consistent with the ordinal role of
$\operatorname{ord}$ in \eqref{eq:graded} and with Rmk.~\ref{rem:st-limits}.
No string-theoretic content is imported; the unique output is the Bell
correlation $\cos2\theta$, not an amplitude.
\end{remark}

% ======================================================================
\section{The transfer operator and the mass gap}\label{sec:gap}
\conditional\ under hypothesis (H1) (README \S3)

\begin{definition}[Bell transfer operator]\label{def:transfer}
Let $T$ be the transfer operator whose integral kernel is the Bell two-point function
derived from~\eqref{I-eq:bell} on the (continuous, hyperfinite) space of measurement
settings. Under hypothesis~(H1) below, $T$ is identified with the Euclidean Yang--Mills
transfer operator.
\end{definition}

\paragraph{Hypothesis (H1).} The Bell two-point function coincides with the Euclidean
Yang--Mills two-point function on the setting space. All results in this section are stated
under~(H1); we do not claim~(H1) itself.

\begin{theorem}[Strict positivity; Jentzsch--Perron]\label{thm:jp}
The kernel of $T$ is strictly positive; consequently its largest eigenvalue is simple and
its eigenvector strictly positive, and $T\ge 1>0$.
\end{theorem}

\begin{proof}
The kernel built from $|\Cmat|=|\cos2\theta|/(2N)$ is positive on the near-standard domain
(Proposition~\ref{prop:exclude} removing the measure-zero singular set). The
Jentzsch--Perron theorem---the integral-operator analogue of Perron--Frobenius---then gives
a simple, isolated largest eigenvalue with positive eigenfunction. The transfer principle
of nonstandard analysis transports this from the hyperfinite to the standard setting.
\end{proof}

\begin{theorem}[Mass gap positivity]\label{thm:gap}
\conditional\ Under hypothesis~(H1), the spectral gap of $T$ yields a positive mass gap
\begin{equation}
\Delta=\frac{\hbar c}{\xi}=\LamQCD>0 .
\end{equation}
The positivity $\Delta>0$ is the mathematical content; the numerical value $\LamQCD$ enters
only through the empirical correlation length $\xi$. The argument is lattice-independent,
living on the continuous/hyperfinite setting space rather than a spacetime lattice.
\end{theorem}

\begin{proof}
By Theorem~\ref{thm:jp} the largest eigenvalue $\lambda_0$ is simple and separated from the
remainder of the spectrum by a finite gap; the mass gap is $\Delta=-\log(\lambda_1/\lambda_0)
\cdot(\hbar c/\xi)>0$ where $\lambda_1<\lambda_0$ is the next eigenvalue. Strict positivity
of the spectral gap is the mathematical assertion; identifying the prefactor with $\LamQCD$
uses the empirical $\xi$ and~(H1).
\end{proof}

\begin{remark}[What is and is not claimed]\label{rem:gap-claim}
Only the positivity $\Delta>0$ is derived. The value $\LamQCD\approx200$~MeV is an
empirical input, not a prediction, and the numerical factor $\approx7.5$ between $\LamQCD$
and the lightest glueball mass remains open.
\end{remark}

% ======================================================================
\section{$SU(3)$ gauge structure}\label{sec:su3}
\conditional

From the colour-singlet qutrit correlation $\Eang=\cos2\theta/6$ (i.e.~\eqref{I-eq:bell}
with $N=3$) one is led towards a Wilson action for an $SU(3)$ gauge structure. This
direction is recorded as conditional: the precise boundary between what is established and
what is in progress---in particular the derivation of the Wilson action from the
correlation rather than its postulation---requires further work and is deferred to the
notebooks. \deferred

% ======================================================================
\section{The quantum correlation: upstream analysis of the Tsirelson bound}\label{sec:qcorr}
\established\ (notebooks 30--33; foundational layer in \texorpdfstring{${}^{*}\RR$}{*R})

The climbing map of \S\ref{sec:climb} reaches $\cos2\theta$ through postulate A2, the
saturation of the Tsirelson bound $2\sqrt2$. We analyse here how much of A2 is itself derived
from the given correlation and how much must be posited, working throughout in nonstandard
analysis: the saturation \emph{value} is asserted on an internal set (the attainment of a
maximum being an internal, hyperfinite statement, not a standard-compactness one), while the
algebraic identities below are first-order and rise to ${}^{*}\RR$ by transfer.

\begin{proposition}[Saturation quality is on-the-given]\label{prop:qual}
The quality that the correlation saturates the quantum bound follows from premises P1--P2
alone. Equivalently: the exclusivity-graph (Cabello--Severini--Winter) quantum bound for CHSH
is the Lov\'asz number $\vartheta=2+\sqrt2$ of the eight-event graph, giving standard
$\text{CHSH}=2\vartheta-4=2\sqrt2$; the orthonormal representation it optimises over requires
exactly the ``inner product $=$ cosine of the relative angle'' geometry that the feature
circle of Part~I, Lemma~\ref{I-lem:bounded} already supplies. The normalised value is
independent of $N$.
\end{proposition}

\begin{proof}
The Lov\'asz number is the optimum of an internal semidefinite program over the feasible set
$\{M\succeq0,\ \operatorname{tr}M=1,\ M_{ij}=0\text{ on edges}\}$~\cite{csw,rabelo}; lifted to ${}^{*}\RR$ this
feasible set is internal and ${}^{*}$-bounded-closed, so by the internal maximum principle the
optimum is attained at an internal point whose standard part is the standard optimiser
(notebook~31, \S31.2; the internality replaces standard compactness). Numerically the optimum
is $2+\sqrt2$. On the lattice $\theta_i=i\varepsilon$ the normalised feature inner product is
$\langle\varphi(\theta_i),\varphi(\theta_j)\rangle\cdot2N=\cos2(\theta_i-\theta_j)$, an identity
holding internally for all $i,j$ (verified symbolically with $\varepsilon$ infinitesimal); this
is precisely the cosine geometry the orthonormal representation demands. The normalisation
removes the amplitude, so the value is $N$-independent.
\end{proof}

\begin{remark}[The bound value is a re-description, not a new principle]\label{rem:redesc}
Proposition~\ref{prop:qual} is not the discovery of an upstream principle selecting $2\sqrt2$;
it is the observation that the saturation \emph{quality} was already contained in P1--P2.
Equivalently, the Lov\'asz $\vartheta$ does not derive the feature angle $2\theta$---it
\emph{presupposes} the cosine geometry that the given circle provides, so ``$\vartheta\to
\cos2\theta\to2\sqrt2$'' is a cycle through a shared geometry, not an upstream-to-downstream
derivation. The value $2\sqrt2$ is the operator-algebraic content of the Tsirelson identity
$B^2=4I-[A_0,A_1][B_0,B_1]$ with $\|[A_0,A_1]\|\le2$~\cite{tsirelson,landau}, the same fact in graph language.
\end{remark}

\begin{remark}[The bound value as an extremal consequence of the $\pm1$ spectrum]\label{rem:extremal-2sqrt2}
The numerical \emph{value} $2\sqrt2$ may be read as the solution of an extremal problem whose
only input is the two-valued spectrum (B3 of Part~I, Principle~\ref{I-prin:bell}). From the
identity $B^2=4I-[A_0,A_1][B_0,B_1]$ and the norm bound
$\|[A_0,A_1]\|\le2\|A_0\|\,\|A_1\|\le2$, valid for any self-adjoint $A_i$ with spectrum in
$[-1,1]$, one gets $\|B\|^2\le4+\|[A_0,A_1]\|\,\|[B_0,B_1]\|\le8$, hence $\|B\|\le2\sqrt2$,
attained on a qubit (verified numerically over random dichotomic observables, notebook~39).
Maximising CHSH subject to the bounded commutator therefore returns $2\sqrt2$ without any
further quantum postulate: the \emph{value} is downstream of the $\pm1$ spectrum. This does
\emph{not} weaken Theorem~\ref{thm:noncomm}: what B3 supplies is the commutator norm
\emph{bound}, whereas the noncommutativity itself---that $[A_0,A_1]\neq0$, the choice to
represent measurements in an associative ${}^{*}$-algebra at all---remains the independent
datum the theorem isolates (a commutative representation would give $\|B\|\le2$, and the
Popescu--Rohrlich box shows non-commutativity is strictly narrower than Bell violation). The
refinement is only that, once the projective $\pm1$ structure is granted, the bound's numerical
value is extremal rather than separately posited; cf.\ the analogous selection-by-extremisation
in the amplitude setting~\cite{cheung2025}.
\end{remark}

\begin{theorem}[Noncommutativity is an independent datum]\label{thm:noncomm}
The noncommutativity that fixes the bound value $2\sqrt2$---the projective, exclusive structure
of measurement---is derived neither from $\cos2\theta$ nor from the standard-part measurement,
and is therefore an external structural input.
\end{theorem}

\begin{proof}
Three exhaustive possibilities are excluded or located. (\emph{$\st$ internal: excluded.}) By
Rmk.~\ref{rem:st-algebra}, $\st$ commutes with commutators and is algebraically neutral, so
noncommutativity is neither created nor destroyed by measurement. (\emph{Derived from
$\cos2\theta$: only halfway.}) By Fine's theorem~\cite{fine}, satisfying all CHSH inequalities is equivalent
to the existence of a single (commutative) joint distribution; since $\cos2\theta$ gives
$\text{CHSH}=2\sqrt2>2$, it entails the \emph{absence} of a commutative model. But absence of a
commutative model does not entail representability by a noncommutative operator algebra: the
Popescu--Rohrlich box~\cite{pr} ($\text{CHSH}=4$, no-signalling) is neither commutative nor quantum, so
noncommutativity (the Tsirelson ceiling $2\sqrt2$) is strictly narrower than mere Bell
violation, and $\cos2\theta$ supplies only the violation. (\emph{Independent datum: located.})
The bound $\|[A_0,A_1]\|\le2$ presupposes self-adjointness, spectrum in $\{\pm1\}$ (two-valued
projective measurement, $A^2=I$), and an associative ${}^{*}$-product---three ingredients absent
from both $\cos2\theta$ and $\st$. Hence noncommutativity is a separate representational datum:
the choice to represent measurements as self-adjoint elements of an associative ${}^{*}$-algebra
(notebook~33, \S\S33.2--33.3).
\end{proof}

\begin{remark}[Necessity is not derivability]\label{rem:nec}
The step that ``a commutative model is impossible'' (Fine) must not be read as ``noncommutativity
is derived from the given.'' Proving a necessary condition---that any reproducing model must fail
to be commutative---does not produce the structure that meets it; the Popescu--Rohrlich box
witnesses a non-commutative-yet-non-quantum alternative. This is the proof-theoretic core of
Theorem~\ref{thm:noncomm}: the bound $2\sqrt2$ rests on a datum (the ${}^{*}$-algebra
representation) that the principle posits rather than derives.
\end{remark}

% ======================================================================
\section{Lorentz invariance: a three-layer emergence}\label{sec:lorentz}
\established\ (notebooks 34--35; foundational layer in \texorpdfstring{${}^{*}\RR$}{*R})

The metric-level $SO(3)$ of Part~I, \S\ref{I-sec:so3} is the spatial rotation group only.
We analyse whether the full Lorentz group $SO(3,1)$ emerges, and find a three-layer
structure in which time (the covering coordinate, an ordered $\RR$) and space (the MDS of
$|\Cmat|$, a metric $\Tcirc^3$) enter on different footings.

\begin{proposition}[Time carries no metric from the static correlation]\label{prop:no-time-metric}
The covering coordinate is totally ordered but, from the static correlation~\eqref{I-eq:bell}
alone, carries no metric: $-\log|\Eang(t)|$ is periodic, hence not a monotone time
distance. A monotone time distance requires a time evolution; if that evolution is a
one-parameter group $U(\tau)=e^{-iH\tau}$ (Stone), the decay is $e^{-\Gamma\tau}$ and the
time distance is the linear $\Gamma\tau$.
\end{proposition}

\begin{theorem}[Three layers of Lorentz emergence]\label{thm:lorentz}
\begin{enumerate}
\item[\rm(L1)] \emph{(Time evolution group $\Rightarrow$ Galilei.)} Adjoining the existence
of a time-evolution group (Stone's $H$) to~\eqref{I-eq:bell} yields the Galilei group:
absolute time (the total order on the covering $\RR$), $SO(3)$, and \emph{commuting}
boosts $t'=t,\ x'=x-vt$. Galilean boosts do not mix time and space, so the mismatch
between the linear time distance and the quadratic spatial distance
($-\log|\cos2\varphi|\sim2\varphi^2$) does not obstruct closure. This adjunction posits no
$\hbar$, $c$, or spectrum, and is the minimal admissible dynamical datum.
\item[\rm(L2)] \emph{(Invariant speed $c$ $\Rightarrow$ local Lorentz.)} Adjoining a
universal evolution rate, i.e.\ an invariant speed $c$, normalises tangent-space
coordinates $T=c\,\Delta\tau$, $X=\sqrt2\,\Delta\varphi$ so that $ds^2=-T^2+X^2$ is boost
invariant and $\mathfrak{so}(3,1)$ closes ($[K_i,K_j]=-\varepsilon_{ijk}J_k$). Local
Lorentz invariance holds in the tangent space.
\item[\rm(L3)] \emph{(Spatial decompactification $\Rightarrow$ global Lorentz.)} Global
boosts are obstructed by the compactness of $\Tcirc^3$ (a hyperbolic orbit leaves any
fundamental domain); global Lorentz invariance requires decompactifying space
($\Tcirc^3\to\RR^3$, a spatial analogue of covering~II that breaks premise P2).
\end{enumerate}
\end{theorem}

\begin{proof}
(L1) The Galilei algebra closes with $[G_i,G_j]=0$, $[G_i,H]=P_i$, $[J_i,G_j]=\varepsilon
G_k$; commuting boosts never mix the time and space sectors, so the differing scaling
dimensions are irrelevant (notebook~35, \S35.1). (L2) In normal coordinates both
directions are linear; the boost $T'=T\cosh w-X\sinh w$, $X'=-T\sinh w+X\cosh w$ leaves
$-T^2+X^2$ invariant, and the six generators of $\mathfrak{so}(3,1)$ satisfy the standard
brackets, verified by matrix computation (notebook~35, \S35.2). (L3) A hyperbolic rotation
is unbounded, incompatible with the compact $\Tcirc^3$; only $\RR^3$ admits global boosts
(notebook~35, \S35.3). The invariant speed $c$ is, under stance~X of Part~I,
Rmk.~\ref{I-rem:c-status}, an external input; under stance~Y it is the empirical datum~(B2).
\end{proof}

\begin{remark}[The correlation is intrinsically non-relativistic]\label{rem:nonrel}
Theorem~\ref{thm:lorentz}(L1) shows that $\cos2\theta$ plus the mere existence of dynamics
yields a \emph{Galilean} (absolute-time, non-relativistic) spacetime. Relativity appears
only at L2, with the invariant speed $c$; under stance~Y this $c$ is supplied empirically
by the spacelike separation of the Bell experiment (Part~I, Rmk.~\ref{I-rem:c-status}),
not derived from the numerical core $\cos2\theta$. The asymmetry between time (order) and
space (metric) of Part~I, Prop.~\ref{I-prop:sep} is what separates Galilei from Lorentz at
the level of the symmetry group.
\end{remark}

% ======================================================================
\section{Climbing map from a minimal principle}\label{sec:climb}
\established\ (notebook 21)

\begin{proposition}[Three additional postulates]\label{prop:climb}
From the minimal principle
\[
\textbf{P0:}\quad\text{there exists a Bell-inequality-violating correlation } C(a,b),
\]
the full form $\Eang=\cos2\theta/(2N)$ is reached by three additional postulates:
\textbf{A1} (homogeneity: $C$ depends only on the relative angle $a-b$), \textbf{A2}
(quantum correlation: saturation of the Tsirelson bound $2\sqrt2$, equivalently a Hilbert
space structure), and \textbf{A3} ($N=3$ QCD identification). Each step is checked by a
CHSH computation.
\end{proposition}

\begin{proof}[Justification]
P0 fixes only that $C$ is a setting-pair function with a CHSH value exceeding $2$;
inhomogeneous correlations can still violate CHSH, so A1 is a genuine added postulate.
Among homogeneous correlations, saturating the Tsirelson bound $|S|=2\sqrt2$ singles out
$E=\cos(\text{relative angle})$, giving A2; the doubling to $\cos2\theta$ is the spin
projective (two-valued) structure carried by A2. Finally A3 fixes the amplitude $1/(2N)$
with $N=3$. The CHSH values at each stage are verified numerically (notebook~21). The
internal structure of A2 is dissected in \S\ref{sec:qcorr}: its saturation \emph{quality} is
on-the-given (Prop.~\ref{prop:qual}), whereas its bound \emph{value} $2\sqrt2$ and the
underlying noncommutativity are external (Thm.~\ref{thm:noncomm}); A2 is thus not a single
indivisible postulate but a quality that is derived atop a posited algebraic datum.
\end{proof}

\begin{remark}[Two independent axes]\label{rem:axes}
The ascent to $\cos2\theta$ (call it axis~I) and the dimension/topology of physical space
(axis~II) are independent: dimension three appears at no stage of axis~I (notebook~21).
This is the Part~II statement of the orthogonality recorded in Part~I,
Rmk.~\ref{I-rem:orthogonal}: $\cos2\theta$ can be climbed without ever fixing the spatial
dimension, which must therefore be supplied separately.
\end{remark}

% ======================================================================
\section{Summary and open problems}\label{sec:math-summary}
The mathematical backbone comprises: the standard-part measurement theory
(\S\ref{sec:st}), under which Bell measurement is a single $\st$ operation that
automatically excludes non-near-standard quantities yet fixes neither dimension, topology,
commutation relations, nor the ordering of spectral modes
(Rmks.~\ref{rem:st-limits}, \ref{rem:st-algebra}, \ref{rem:st-ordering}); the spectral/MDS
recovery of geometry (\S\ref{sec:mds}), giving $\Tcirc^{3}$ with orthogonal axes; the
upstream analysis of the quantum correlation (\S\ref{sec:qcorr}), separating the saturation
quality (on-the-given, Prop.~\ref{prop:qual}) from the bound value $2\sqrt2$ and the
noncommutativity that fixes it (external, Thm.~\ref{thm:noncomm}); and the climbing map
(\S\ref{sec:climb}), reducing $\cos2\theta$ to the postulates A1--A3. The mass-gap positivity
(\S\ref{sec:gap}) and the $SU(3)$ Wilson structure (\S\ref{sec:su3}) are conditional, the
former on hypothesis~(H1).

The structural question of axis~I---whether the quantum-correlation postulate A2 descends
from a more minimal principle---is now resolved in outline. The saturation \emph{quality}
does descend (it is the constant-norm circle of P1--P2 re-described in graph-theoretic terms,
Prop.~\ref{prop:qual} and Rmk.~\ref{rem:redesc}); the bound \emph{value} $2\sqrt2$ descends
only partway---given the projective $\pm1$ spectrum (B3) it is the extremum of CHSH under the
commutator bound $\|[A_0,A_1]\|\le2$ (Rmk.~\ref{rem:extremal-2sqrt2}), but the
noncommutativity that makes that commutator nonzero in the first place remains an independent
representational datum (Thm.~\ref{thm:noncomm}). The open problems are therefore the physical
justification of the external inputs rather than their further reduction: the upstream origin
of spatial dimension three, independent of $\cos2\theta$ (axis~II of Rmk.~\ref{rem:axes}); the
origin of the projective/noncommutative measurement structure that sets $2\sqrt2$; the $N=3$
identification, to be addressed via the $SU(3)$ Wilson action (\S\ref{sec:su3}); the
order assignment of spectral modes, which the rank leaves undetermined
(Rmk.~\ref{rem:st-ordering}) and which controls the conditional bridge to the companion
paper; and the
numerical factor $\approx7.5$ between $\LamQCD$ and the lightest glueball mass.

% ======================================================================
\begin{thebibliography}{9}
\bibitem{albeverio} S. Albeverio et al., \emph{Nonstandard Methods in Stochastic Analysis and Mathematical Physics}, Academic Press (1986).
\bibitem{yamashita} H. Yamashita, \emph{Nonstandard methods in quantum field theory I: a hyperfinite formalism of scalar fields}, Int. J. Theor. Phys. \textbf{41}, 511 (2002).
\bibitem{nsa-qm} \emph{A realistic model for completing quantum mechanics} (nonstandard hyperfinite real axis), arXiv:2104.12701 (2021).
\bibitem{finberg} J. S. Finberg, \emph{Bell nonlocality as a covariance obstruction in locally covariant quantum field theory}, arXiv:2512.18603 (2025).
\bibitem{karlsson} A. Karlsson, \emph{Space-time emergence from individual interactions}, arXiv:1806.05710 (2018).
\bibitem{csw} A. Cabello, S. Severini, and A. Winter, \emph{Graph-theoretic approach to quantum correlations}, Phys. Rev. Lett. \textbf{112}, 040401 (2014).
\bibitem{rabelo} R. Rabelo, C. Duarte, A. J. L\'opez-Tarrida, M. T. Cunha, and A. Cabello, \emph{Multigraph approach to quantum non-locality}, J. Phys. A \textbf{47}, 424021 (2014).
\bibitem{tsirelson} B. S. Cirel'son (Tsirelson), \emph{Quantum generalizations of Bell's inequality}, Lett. Math. Phys. \textbf{4}, 93 (1980).
\bibitem{landau} L. J. Landau, \emph{On the violation of Bell's inequality in quantum theory}, Phys. Lett. A \textbf{120}, 54 (1987).
\bibitem{fine} A. Fine, \emph{Hidden variables, joint probability, and the Bell inequalities}, Phys. Rev. Lett. \textbf{48}, 291 (1982).
\bibitem{pr} S. Popescu and D. Rohrlich, \emph{Quantum nonlocality as an axiom}, Found. Phys. \textbf{24}, 379 (1994).
\bibitem{emergentdist2025} B. Leighton-Trudel, \emph{Emergent distance and metricity of mutual information in 1D quantum chains}, arXiv:2507.09749 (2025).
\bibitem{persistenthom} B. Olsthoorn, \emph{Persistent homology of quantum entanglement}, Phys. Rev. B \textbf{107}, 115174 (2023); arXiv:2110.10214 (2021).
\bibitem{rank2paper} \emph{Rank-two rigidity of Bell correlations and the bipartite character of emergent geometry}, companion paper (this project), \S5.5--5.6 (inverse problem, harmonic tower, NSA bridge).
\bibitem{cheung2025} C. Cheung, G. N. Remmen, F. Sciotti, M. Tarquini, \emph{Strings from almost nothing}, arXiv:2508.09246 (2025).
\end{thebibliography}

\end{document}
